\documentclass[12pt,a4paper]{article}

% Codificación y lenguaje
\usepackage[utf8]{inputenc}
\usepackage[T1]{fontenc}
\usepackage[spanish,es-nodecimaldot]{babel}

% Matemáticas
\usepackage{amsmath,amssymb,amsthm}
\usepackage{cancel}

% Diseño
\usepackage[margin=2.5cm]{geometry}
\usepackage{multicol}
\usepackage{booktabs}
\usepackage{array}
\usepackage{enumitem}

% Gráficos
\usepackage{graphicx}
\usepackage{tikz}
\usepackage{pgfplots}
\pgfplotsset{compat=1.18}
\usepgfplotslibrary{fillbetween} % por buena práctica

% Hiperenlaces
\usepackage{hyperref} % antes estaba hyperref (mal escrito)
\hypersetup{
    colorlinks=true,
    linkcolor=blue,
    urlcolor=blue
}

% Cajas
\usepackage{tcolorbox}
\tcbuselibrary{skins,breakable}

\newtcolorbox{funcbox}[2][]{
    colback=blue!5!white,
    colframe=blue!75!black,
    fonttitle=\bfseries,
    title={#2},
    #1
}

% Definición de dominio y recorrido en español
\newcommand{\Dominio}{\ensuremath{\mathrm{Dom}}}
\newcommand{\Recorrido}{\ensuremath{\mathrm{Rec}}}

\begin{document}

% Portada
\title{\textbf{Funciones Elementales y sus Propiedades}\\ \Large Un repaso completo para estudiantes}
\author{}
\date{\today}
\maketitle

\tableofcontents
\newpage

%=====================================================================
% APARTADO 1: FUNCIONES INICIALES
%=====================================================================
\section{Funciones Iniciales que Deben Conocer los Estudiantes}

A continuación presentamos las funciones fundamentales, sus gráficas, dominio y recorrido (rango). Se incluyen tanto las formas canónicas como variantes con desplazamientos.

%--- Lineal ---
\subsection{Función Lineal}
\begin{funcbox}{Función Lineal $f(x)=mx+b$}
\begin{minipage}[t]{0.55\textwidth}
    \textbf{Definición:} $f:\mathbb{R}\to\mathbb{R}$, $f(x)=mx+b$, con $m,b\in\mathbb{R}$.\\
    \textbf{Dominio:} $\Dominio(f)=\mathbb{R}$\\
    \textbf{Recorrido:} 
    \begin{itemize}
        \item Si $m\neq 0$: $\Recorrido(f)=\mathbb{R}$
        \item Si $m=0$: función constante, $\Recorrido(f)=\{b\}$
    \end{itemize}
    \textbf{Características:}
    \begin{itemize}
        \item $m$ es la pendiente, indica la inclinación.
        \item $b$ es la ordenada al origen (corte con eje $y$).
        \item Gráfica: línea recta.
        \item Si $m>0$, la función es creciente; si $m<0$, decreciente.
    \end{itemize}
    \textbf{Ejemplo gráfico:} $f(x)=2x-1$
\end{minipage}\hfill
\begin{minipage}[t]{0.4\textwidth}
    \centering
    \begin{tikzpicture}
        \begin{axis}[
            axis lines=middle,
            xlabel=$x$, ylabel=$y$,
            xmin=-2.5, xmax=2.5,
            ymin=-3.5, ymax=3.5,
            grid=both,
            minor tick num=1,
            width=6cm, height=5cm,
            title={$f(x)=2x-1$}
        ]
        \addplot[domain=-2:2, samples=100, blue, thick] {2*x-1};
        \end{axis}
    \end{tikzpicture}
\end{minipage}
\end{funcbox}

% --- Lineal con desplazamiento ---
\subsection{Función Lineal con Desplazamiento}
\begin{funcbox}{Función Lineal Desplazada: $f(x)=m(x-h)+k$}
\begin{minipage}[t]{0.55\textwidth}
    \textbf{Forma general:} $f(x)=m(x-h)+k$\\
    Pasa por el punto $(h,k)$ y tiene pendiente $m$.\\
    \textbf{Dominio:} $\mathbb{R}$\\
    \textbf{Recorrido:} $\mathbb{R}$ (si $m\neq0$).\\
    \textbf{Características:}
    \begin{itemize}
        \item Es la recta con pendiente $m$ que pasa por $(h,k)$.
        \item Si $h=0$, se obtiene la forma $mx+k$; si $k=0$ es $m(x-h)$.
    \end{itemize}
    \textbf{Ejemplo:} $f(x)=2(x-1)+3$ (pasa por $(1,3)$, pendiente $2$)
\end{minipage}\hfill
\begin{minipage}[t]{0.4\textwidth}
    \centering
    \begin{tikzpicture}
        \begin{axis}[
            axis lines=middle,
            xlabel=$x$, ylabel=$y$,
            xmin=-1, xmax=3,
            ymin=-1, ymax=6,
            grid=both,
            minor tick num=1,
            width=6cm, height=5cm,
            title={$f(x)=2(x-1)+3$}
        ]
        \addplot[domain=-0.5:2.5, samples=100, blue, thick] {2*(x-1)+3};
        \addplot[mark=*, only marks, red] coordinates {(1,3)};
        \end{axis}
    \end{tikzpicture}
\end{minipage}
\end{funcbox}

%--- Cuadrática ---
\subsection{Función Cuadrática}
\begin{funcbox}{Función Cuadrática $f(x)=ax^2+bx+c$ \quad ($a\neq 0$)}
\begin{minipage}[t]{0.55\textwidth}
    \textbf{Forma canónica:} $f(x)=a(x-h)^2+k$ con vértice $(h,k)$.\\
    \textbf{Dominio:} $\Dominio(f)=\mathbb{R}$\\
    \textbf{Recorrido:}
    \begin{itemize}
        \item Si $a>0$: $[k,\infty)$
        \item Si $a<0$: $(-\infty, k]$
    \end{itemize}
    \textbf{Características:}
    \begin{itemize}
        \item Gráfica: parábola vertical.
        \item $a$ determina la abertura: $|a|>1$, más estrecha; $0<|a|<1$, más ancha.
        \item El vértice es el punto mínimo (si $a>0$) o máximo (si $a<0$).
        \item Eje de simetría: recta $x=h$.
    \end{itemize}
    \textbf{Ejemplo:} $f(x)=x^2$ (vértice $(0,0)$, $a=1$)
\end{minipage}\hfill
\begin{minipage}[t]{0.4\textwidth}
    \centering
    \begin{tikzpicture}
        \begin{axis}[
            axis lines=middle,
            xlabel=$x$, ylabel=$y$,
            xmin=-2.5, xmax=2.5,
            ymin=-1, ymax=4,
            grid=both,
            minor tick num=1,
            width=6cm, height=5cm,
            title={$f(x)=x^2$}
        ]
        \addplot[domain=-2:2, samples=100, blue, thick] {x^2};
        \end{axis}
    \end{tikzpicture}
\end{minipage}
\end{funcbox}

% --- Cuadrática desplazada ---
\subsection{Función Cuadrática Desplazada}
\begin{funcbox}{Función Cuadrática Desplazada: $f(x)=a(x-h)^2+k$}
\begin{minipage}[t]{0.55\textwidth}
    \textbf{Ejemplo concreto:} $f(x)=2(x-1)^2+3$\\
    Vértice: $(1,3)$; $a=2>0$, parábola hacia arriba.\\
    \textbf{Dominio:} $\mathbb{R}$\\
    \textbf{Recorrido:} $[3,\infty)$\\
    \textbf{Características:}
    \begin{itemize}
        \item El vértice se ha trasladado a $(h,k)$.
        \item El eje de simetría es $x=h$.
        \item El signo de $a$ sigue determinando la orientación.
    \end{itemize}
\end{minipage}\hfill
\begin{minipage}[t]{0.4\textwidth}
    \centering
    \begin{tikzpicture}
        \begin{axis}[
            axis lines=middle,
            xlabel=$x$, ylabel=$y$,
            xmin=-1, xmax=3.5,
            ymin=-2, ymax=8,
            grid=both,
            minor tick num=1,
            width=6cm, height=5cm,
            title={$f(x)=2(x-1)^2+3$}
        ]
        \addplot[domain=-0.5:2.5, samples=100, blue, thick] {2*(x-1)^2+3};
        \addplot[mark=*, only marks, red] coordinates {(1,3)};
        \end{axis}
    \end{tikzpicture}
\end{minipage}
\end{funcbox}

%--- Valor absoluto ---
\subsection{Función Valor Absoluto}
\begin{funcbox}{Función Valor Absoluto $f(x)=|x|$}
\begin{minipage}[t]{0.55\textwidth}
    \textbf{Definición:} $f(x)=|x|=\begin{cases} 
                x, & x\geq 0 \\
                -x, & x<0 
            \end{cases}$\\
    \textbf{Dominio:} $\Dominio(f)=\mathbb{R}$\\
    \textbf{Recorrido:} $[0,\infty)$\\
    \textbf{Características:}
    \begin{itemize}
        \item Gráfica en forma de ``V''.
        \item Punto anguloso en el origen.
        \item Función par: $f(-x)=f(x)$.
        \item Decreciente en $(-\infty,0]$, creciente en $[0,\infty)$.
    \end{itemize}
    \textbf{Gráfica:}
\end{minipage}\hfill
\begin{minipage}[t]{0.4\textwidth}
    \centering
    \begin{tikzpicture}
        \begin{axis}[
            axis lines=middle,
            xlabel=$x$, ylabel=$y$,
            xmin=-2.5, xmax=2.5,
            ymin=-0.5, ymax=3.5,
            grid=both,
            minor tick num=1,
            width=6cm, height=5cm,
            title={$f(x)=|x|$}
        ]
        \addplot[domain=-2.5:2.5, samples=200, blue, thick] {abs(x)};
        \end{axis}
    \end{tikzpicture}
\end{minipage}
\end{funcbox}

% --- Valor absoluto con desplazamiento ---
\subsection{Función Valor Absoluto con Desplazamiento}
\begin{funcbox}{Función Valor Absoluto Desplazada: $f(x)=a|x-h|+k$}
\begin{minipage}[t]{0.55\textwidth}
    \textbf{Forma general:} $f(x)=a|x-h|+k$\\
    Vértice en $(h,k)$.\\
    \textbf{Dominio:} $\mathbb{R}$\\
    \textbf{Recorrido:}
    \begin{itemize}
        \item Si $a>0$: $[k,\infty)$
        \item Si $a<0$: $(-\infty,k]$
    \end{itemize}
    \textbf{Ejemplo:} $f(x)=|x-2|+1$ (vértice $(2,1)$, abre hacia arriba)
\end{minipage}\hfill
\begin{minipage}[t]{0.4\textwidth}
    \centering
    \begin{tikzpicture}
        \begin{axis}[
            axis lines=middle,
            xlabel=$x$, ylabel=$y$,
            xmin=-1, xmax=4.5,
            ymin=-1, ymax=5,
            grid=both,
            minor tick num=1,
            width=6cm, height=5cm,
            title={$f(x)=|x-2|+1$}
        ]
        \addplot[domain=-0.5:4.5, samples=200, blue, thick] {abs(x-2)+1};
        \addplot[mark=*, only marks, red] coordinates {(2,1)};
        \end{axis}
    \end{tikzpicture}
\end{minipage}
\end{funcbox}

%--- Cúbica (sin desplazamiento) ---
\subsection{Función Cúbica (forma básica)}
\begin{funcbox}{Función Cúbica $f(x)=x^3$}
\begin{minipage}[t]{0.55\textwidth}
    \textbf{Definición:} $f(x)=x^3$\\
    \textbf{Dominio:} $\mathbb{R}$\\
    \textbf{Recorrido:} $\mathbb{R}$\\
    \textbf{Características:}
    \begin{itemize}
        \item Gráfica en forma de ``S'' alargada.
        \item Función impar: $f(-x)=-f(x)$ (simetría respecto al origen).
        \item Creciente en todo $\mathbb{R}$.
        \item Punto de inflexión en $(0,0)$.
    \end{itemize}
\end{minipage}\hfill
\begin{minipage}[t]{0.4\textwidth}
    \centering
    \begin{tikzpicture}
        \begin{axis}[
            axis lines=middle,
            xlabel=$x$, ylabel=$y$,
            xmin=-2.5, xmax=2.5,
            ymin=-4, ymax=4,
            grid=both,
            minor tick num=1,
            width=6cm, height=5cm,
            title={$f(x)=x^3$}
        ]
        \addplot[domain=-2:2, samples=200, blue, thick] {x^3};
        \end{axis}
    \end{tikzpicture}
\end{minipage}
\end{funcbox}

%--- Cúbica con desplazamiento ---
\subsection{Función Cúbica con Desplazamiento}
\begin{funcbox}{Función Cúbica desplazada: $f(x)=a(x-h)^3+k$}
\begin{minipage}[t]{0.55\textwidth}
    \textbf{Forma general:} $f(x)=a(x-h)^3+k$.\\
    \textbf{Dominio:} $\mathbb{R}$\\
    \textbf{Recorrido:} $\mathbb{R}$ (si $a\neq 0$).\\
    \textbf{Características:}
    \begin{itemize}
        \item El punto de inflexión se traslada a $(h,k)$.
        \item La gráfica mantiene la forma cúbica básica pero desplazada.
        \item Si $a<0$, se refleja verticalmente (decreciente global).
    \end{itemize}
    \textbf{Ejemplo:} $f(x)=(x-1)^3+2$ (desplazamiento a la derecha 1, arriba 2)
\end{minipage}\hfill
\begin{minipage}[t]{0.4\textwidth}
    \centering
    \begin{tikzpicture}
        \begin{axis}[
            axis lines=middle,
            xlabel=$x$, ylabel=$y$,
            xmin=-1.5, xmax=3.5,
            ymin=-2, ymax=6,
            grid=both,
            minor tick num=1,
            width=6cm, height=5cm,
            title={$f(x)=(x-1)^3+2$}
        ]
        \addplot[domain=-1:3, samples=200, blue, thick] {(x-1)^3+2};
        \end{axis}
    \end{tikzpicture}
\end{minipage}
\end{funcbox}

%--- Raíz cuadrada ---
\subsection{Función Raíz Cuadrada}
\begin{funcbox}{Función Raíz Cuadrada $f(x)=\sqrt{x}$}
\begin{minipage}[t]{0.55\textwidth}
    \textbf{Definición:} $f(x)=\sqrt{x}=x^{1/2}$ para $x\geq 0$.\\
    \textbf{Dominio:} $[0,\infty)$\\
    \textbf{Recorrido:} $[0,\infty)$\\
    \textbf{Características:}
    \begin{itemize}
        \item Creciente en todo su dominio.
        \item Cóncava hacia abajo.
        \item Punto de inicio en $(0,0)$.
        \item La pendiente decrece al aumentar $x$.
    \end{itemize}
\end{minipage}\hfill
\begin{minipage}[t]{0.4\textwidth}
    \centering
    \begin{tikzpicture}
        \begin{axis}[
            axis lines=middle,
            xlabel=$x$, ylabel=$y$,
            xmin=-0.5, xmax=5,
            ymin=-0.5, ymax=3,
            grid=both,
            minor tick num=1,
            width=6cm, height=5cm,
            title={$f(x)=\sqrt{x}$}
        ]
        \addplot[domain=0:5, samples=200, blue, thick] {sqrt(x)};
        \end{axis}
    \end{tikzpicture}
\end{minipage}
\end{funcbox}

% --- Raíz cuadrada con desplazamiento ---
\subsection{Función Raíz Cuadrada con Desplazamiento}
\begin{funcbox}{Función Raíz Cuadrada Desplazada: $f(x)=a\sqrt{x-h}+k$}
\begin{minipage}[t]{0.55\textwidth}
    \textbf{Forma general:} $f(x)=a\sqrt{x-h}+k$ con $x\ge h$.\\
    Punto de inicio: $(h,k)$.\\
    \textbf{Dominio:} $[h,\infty)$\\
    \textbf{Recorrido:}
    \begin{itemize}
        \item Si $a>0$: $[k,\infty)$
        \item Si $a<0$: $(-\infty,k]$
    \end{itemize}
    \textbf{Ejemplo:} $f(x)=\sqrt{x-2}+1$ (inicio en $(2,1)$)
\end{minipage}\hfill
\begin{minipage}[t]{0.4\textwidth}
    \centering
    \begin{tikzpicture}
        \begin{axis}[
            axis lines=middle,
            xlabel=$x$, ylabel=$y$,
            xmin=1.5, xmax=6,
            ymin=-0.5, ymax=4,
            grid=both,
            minor tick num=1,
            width=6cm, height=5cm,
            title={$f(x)=\sqrt{x-2}+1$}
        ]
        \addplot[domain=2:6, samples=200, blue, thick] {sqrt(x-2)+1};
        \addplot[mark=*, only marks, red] coordinates {(2,1)};
        \end{axis}
    \end{tikzpicture}
\end{minipage}
\end{funcbox}

%--- Raíz cúbica ---
\subsection{Función Raíz Cúbica}
\begin{funcbox}{Función Raíz Cúbica $f(x)=\sqrt[3]{x}=x^{1/3}$}
\begin{minipage}[t]{0.55\textwidth}
    \textbf{Definición:} $f(x)=\sqrt[3]{x}$ para todo $x\in\mathbb{R}$.\\
    \textbf{Dominio:} $\mathbb{R}$\\
    \textbf{Recorrido:} $\mathbb{R}$\\
    \textbf{Características:}
    \begin{itemize}
        \item Función impar (simétrica respecto al origen).
        \item Creciente en todo $\mathbb{R}$.
        \item Punto de inflexión en $(0,0)$.
        \item Para $x>0$, $f(x)>0$; para $x<0$, $f(x)<0$.
    \end{itemize}
\end{minipage}\hfill
\begin{minipage}[t]{0.4\textwidth}
    \centering
    \begin{tikzpicture}
        \begin{axis}[
            axis lines=middle,
            xlabel=$x$, ylabel=$y$,
            xmin=-2.5, xmax=2.5,
            ymin=-2.5, ymax=2.5,
            grid=both,
            minor tick num=1,
            width=6cm, height=5cm,
            title={$f(x)=\sqrt[3]{x}$}
        ]
        \addplot[domain=-2:2, samples=200, blue, thick] {sign(x)*abs(x)^(1/3)};
        \end{axis}
    \end{tikzpicture}
\end{minipage}
\end{funcbox}

% --- Raíz cúbica con desplazamiento ---
\subsection{Función Raíz Cúbica con Desplazamiento}
\begin{funcbox}{Función Raíz Cúbica Desplazada: $f(x)=a\sqrt[3]{x-h}+k$}
\begin{minipage}[t]{0.55\textwidth}
    \textbf{Forma general:} $f(x)=a\sqrt[3]{x-h}+k$\\
    Punto de inflexión: $(h,k)$.\\
    \textbf{Dominio:} $\mathbb{R}$\\
    \textbf{Recorrido:} $\mathbb{R}$\\
    \textbf{Ejemplo:} $f(x)=\frac{1}{2}\sqrt[3]{x+1}+2$ (inflexión en $(-1,2)$)
\end{minipage}\hfill
\begin{minipage}[t]{0.4\textwidth}
    \centering
    \begin{tikzpicture}
        \begin{axis}[
            axis lines=middle,
            xlabel=$x$, ylabel=$y$,
            xmin=-3, xmax=3,
            ymin=0, ymax=4,
            grid=both,
            minor tick num=1,
            width=6cm, height=5cm,
            title={$f(x)=\frac{1}{2}\sqrt[3]{x+1}+2$}
        ]
        \addplot[domain=-3:3, samples=200, blue, thick] {0.5*sign(x+1)*abs(x+1)^(1/3)+2};
        \addplot[mark=*, only marks, red] coordinates {(-1,2)};
        \end{axis}
    \end{tikzpicture}
\end{minipage}
\end{funcbox}

%--- Logarítmica ---
\subsection{Función Logarítmica}
\begin{funcbox}{Función Logarítmica $f(x)=\log_a(x)$ con $a>0$, $a\neq 1$}
\begin{minipage}[t]{0.55\textwidth}
    \textbf{Definición:} $y=\log_a(x) \iff a^y=x$.\\
    \textbf{Dominio:} $(0,\infty)$\\
    \textbf{Recorrido:} $\mathbb{R}$\\
    \textbf{Características:}
    \begin{itemize}
        \item Si $a>1$, es creciente; si $0<a<1$, decreciente.
        \item Asíntota vertical en $x=0$.
        \item Pasa por $(1,0)$ y $(a,1)$.
        \item Logaritmo natural: $\ln(x):=\log_e(x)$.
    \end{itemize}
    \textbf{Ejemplo:} $f(x)=\ln(x)$
\end{minipage}\hfill
\begin{minipage}[t]{0.4\textwidth}
    \centering
    \begin{tikzpicture}
        \begin{axis}[
            axis lines=middle,
            xlabel=$x$, ylabel=$y$,
            xmin=-0.5, xmax=5,
            ymin=-3, ymax=2.5,
            grid=both,
            minor tick num=1,
            width=6cm, height=5cm,
            title={$f(x)=\ln(x)$}
        ]
        \addplot[domain=0.05:5, samples=200, blue, thick] {ln(x)};
        \draw[dashed] (axis cs:0,-3) -- (axis cs:0,2.5);
        \end{axis}
    \end{tikzpicture}
\end{minipage}
\end{funcbox}

% --- Logarítmica con desplazamiento ---
\subsection{Función Logarítmica con Desplazamiento}
\begin{funcbox}{Función Logarítmica Desplazada: $f(x)=a\log_b(x-h)+k$}
\begin{minipage}[t]{0.55\textwidth}
    \textbf{Forma general:} $f(x)=a\log_b(x-h)+k$\\
    Asíntota vertical en $x=h$.\\
    \textbf{Dominio:} $(h,\infty)$\\
    \textbf{Recorrido:} $\mathbb{R}$\\
    \textbf{Ejemplo:} $f(x)=\ln(x-1)+2$ (asíntota $x=1$, desplazada arriba 2)
\end{minipage}\hfill
\begin{minipage}[t]{0.4\textwidth}
    \centering
    \begin{tikzpicture}
        \begin{axis}[
            axis lines=middle,
            xlabel=$x$, ylabel=$y$,
            xmin=0.5, xmax=5,
            ymin=-2, ymax=5,
            grid=both,
            minor tick num=1,
            width=6cm, height=5cm,
            title={$f(x)=\ln(x-1)+2$}
        ]
        \addplot[domain=1.05:5, samples=200, blue, thick] {ln(x-1)+2};
        \draw[dashed] (axis cs:1,-2) -- (axis cs:1,5);
        \end{axis}
    \end{tikzpicture}
\end{minipage}
\end{funcbox}

%--- Exponencial ---
\subsection{Función Exponencial}
\begin{funcbox}{Función Exponencial $f(x)=a^x$ con $a>0$, $a\neq 1$}
\begin{minipage}[t]{0.55\textwidth}
    \textbf{Definición:} $f(x)=a^x$. Caso típico: $f(x)=e^x$.\\
    \textbf{Dominio:} $\mathbb{R}$\\
    \textbf{Recorrido:} $(0,\infty)$\\
    \textbf{Características:}
    \begin{itemize}
        \item Si $a>1$, creciente; si $0<a<1$, decreciente.
        \item Asíntota horizontal en $y=0$.
        \item Pasa por $(0,1)$ y $(1,a)$.
        \item Inversa de la función logarítmica.
    \end{itemize}
    \textbf{Ejemplo:} $f(x)=e^x$
\end{minipage}\hfill
\begin{minipage}[t]{0.4\textwidth}
    \centering
    \begin{tikzpicture}
        \begin{axis}[
            axis lines=middle,
            xlabel=$x$, ylabel=$y$,
            xmin=-3, xmax=3,
            ymin=-1, ymax=5,
            grid=both,
            minor tick num=1,
            width=6cm, height=5cm,
            title={$f(x)=e^x$}
        ]
        \addplot[domain=-3:1.6, samples=200, blue, thick] {exp(x)};
        \draw[dashed] (axis cs:-3,0) -- (axis cs:3,0);
        \end{axis}
    \end{tikzpicture}
\end{minipage}
\end{funcbox}

% --- Exponencial con desplazamiento ---
\subsection{Función Exponencial con Desplazamiento}
\begin{funcbox}{Función Exponencial Desplazada: $f(x)=a\cdot b^{(x-h)}+k$}
\begin{minipage}[t]{0.55\textwidth}
    \textbf{Forma general:} $f(x)=a\cdot b^{(x-h)}+k$\\
    Asíntota horizontal en $y=k$.\\
    \textbf{Dominio:} $\mathbb{R}$\\
    \textbf{Recorrido:}
    \begin{itemize}
        \item Si $a>0$: $(k,\infty)$
        \item Si $a<0$: $(-\infty,k)$
    \end{itemize}
    \textbf{Ejemplo:} $f(x)=2\cdot 3^{(x-1)}-2$ (asíntota $y=-2$, pasa por $(1,0)$)
\end{minipage}\hfill
\begin{minipage}[t]{0.4\textwidth}
    \centering
    \begin{tikzpicture}
        \begin{axis}[
            axis lines=middle,
            xlabel=$x$, ylabel=$y$,
            xmin=-2, xmax=3,
            ymin=-3, ymax=5,
            grid=both,
            minor tick num=1,
            width=6cm, height=5cm,
            title={$f(x)=2\cdot 3^{(x-1)}-2$}
        ]
        \addplot[domain=-2:2.5, samples=200, blue, thick] {2*3^(x-1)-2};
        \draw[dashed] (axis cs:-2,-2) -- (axis cs:3,-2);
        \end{axis}
    \end{tikzpicture}
\end{minipage}
\end{funcbox}

%--- Racional ---
\subsection{Función Racional de la forma $y=\dfrac{ax+b}{cx+d}$}
\begin{funcbox}{Función Homográfica $f(x)=\dfrac{ax+b}{cx+d}$ \quad ($c\neq 0$, $ad\neq bc$)}
\begin{minipage}[t]{0.55\textwidth}
    \textbf{Definición:} Cociente de polinomios lineales.\\
    \textbf{Dominio:} $\mathbb{R}\setminus\left\{-\dfrac{d}{c}\right\}$\\
    \textbf{Recorrido:} $\mathbb{R}\setminus\left\{\dfrac{a}{c}\right\}$ (asíntota horizontal)\\
    \textbf{Características:}
    \begin{itemize}
        \item Asíntota vertical: $x=-\dfrac{d}{c}$.
        \item Asíntota horizontal: $y=\dfrac{a}{c}$.
        \item Hipérbola equilátera trasladada.
        \item Monótona por ramas.
    \end{itemize}
    \textbf{Ejemplo:} $f(x)=\dfrac{2x+1}{x-3}$; Dom: $\mathbb{R}\setminus\{3\}$, Rec: $\mathbb{R}\setminus\{2\}$.
\end{minipage}\hfill
\begin{minipage}[t]{0.4\textwidth}
    \centering
    \begin{tikzpicture}
        \begin{axis}[
            axis lines=middle,
            xlabel=$x$, ylabel=$y$,
            xmin=-2, xmax=8,
            ymin=-5, ymax=7,
            grid=both,
            minor tick num=1,
            width=6cm, height=5cm,
            title={$f(x)=\frac{2x+1}{x-3}$},
            restrict y to domain=-5:7
        ]
        \addplot[domain=-2:2.8, samples=200, blue, thick] {(2*x+1)/(x-3)};
        \addplot[domain=3.2:8, samples=200, blue, thick] {(2*x+1)/(x-3)};
        \draw[dashed] (axis cs:3,-5) -- (axis cs:3,7);
        \draw[dashed] (axis cs:-2,2) -- (axis cs:8,2);
        \end{axis}
    \end{tikzpicture}
\end{minipage}
\end{funcbox}

% --- Proporcionalidad inversa desplazada ---
\subsection{Función Proporcionalidad Inversa Desplazada}
\begin{funcbox}{Función $f(x)=\dfrac{a}{x-h}+k$}
\begin{minipage}[t]{0.55\textwidth}
    \textbf{Forma general:} $f(x)=\dfrac{a}{x-h}+k$\\
    Es la hipérbola $y=1/x$ trasladada.\\
    \textbf{Dominio:} $\mathbb{R}\setminus\{h\}$\\
    \textbf{Recorrido:} $\mathbb{R}\setminus\{k\}$\\
    Asíntotas: $x=h$ (vertical) e $y=k$ (horizontal).\\
    \textbf{Ejemplo:} $f(x)=\dfrac{2}{x-1}+3$ (asíntotas $x=1$, $y=3$)
\end{minipage}\hfill
\begin{minipage}[t]{0.4\textwidth}
    \centering
    \begin{tikzpicture}
        \begin{axis}[
            axis lines=middle,
            xlabel=$x$, ylabel=$y$,
            xmin=-2, xmax=4,
            ymin=-2, ymax=8,
            grid=both,
            minor tick num=1,
            width=6cm, height=5cm,
            title={$f(x)=\frac{2}{x-1}+3$},
            restrict y to domain=-2:8
        ]
        \addplot[domain=-2:0.85, samples=200, blue, thick] {2/(x-1)+3};
        \addplot[domain=1.15:4, samples=200, blue, thick] {2/(x-1)+3};
        \draw[dashed] (axis cs:1,-2) -- (axis cs:1,8);
        \draw[dashed] (axis cs:-2,3) -- (axis cs:4,3);
        \end{axis}
    \end{tikzpicture}
\end{minipage}
\end{funcbox}

%--- Por tramos ---
\subsection{Función Definida por Tramos}
\begin{funcbox}{Función a trozos}
\begin{minipage}[t]{0.55\textwidth}
    \textbf{Definición:} La función se define con expresiones diferentes en intervalos disjuntos.\\
    \textbf{Dominio:} Unión de los intervalos de definición.\\
    \textbf{Recorrido:} Se determina analizando cada tramo.\\
    \textbf{Ejemplo:} 
    $f(x)=\begin{cases}
        x^2, & x<0 \\
        x+1, & x\geq 0
    \end{cases}$\\
    Dominio: $\mathbb{R}$; Recorrido: $(0,\infty)\cup[1,\infty) = (0,\infty)$ (porque $x^2>0$ para $x<0$, y $x+1\geq1$ para $x\geq0$).\\
    Puntos de ruptura pueden presentar saltos o cambios de fórmula.
\end{minipage}\hfill
\begin{minipage}[t]{0.4\textwidth}
    \centering
    \begin{tikzpicture}
        \begin{axis}[
            axis lines=middle,
            xlabel=$x$, ylabel=$y$,
            xmin=-3, xmax=3,
            ymin=-1, ymax=5,
            grid=both,
            minor tick num=1,
            width=6cm, height=5cm,
            title={Por tramos}
        ]
        \addplot[domain=-3:0, samples=200, blue, thick, solid] {x^2};
        \addplot[domain=0:3, samples=200, blue, thick, solid] {x+1};
        % Círculo abierto en (0,0) (no pertenece a la primera rama)
        \addplot[mark=*, fill=white, only marks, blue] coordinates {(0,0)};
        % Punto cerrado en (0,1) (sí pertenece a la segunda rama)
        \addplot[mark=*, only marks, blue] coordinates {(0,1)};
        \end{axis}
    \end{tikzpicture}
\end{minipage}
\end{funcbox}

%=====================================================================
% APARTADO 2: PROPIEDADES DE LOGARITMOS
%=====================================================================
\section{Propiedades de los Logaritmos}

Sea $a>0$, $a\neq 1$, y $x,y>0$, $n\in\mathbb{R}$. Se cumple:

\begin{tcolorbox}[colback=green!5, colframe=green!70!black, title=Propiedades Fundamentales]
\begin{enumerate}[label=\textbf{\arabic*.}]
    \item \textbf{Logaritmo del producto:} $\log_a(xy) = \log_a(x) + \log_a(y)$.
    \item \textbf{Logaritmo del cociente:} $\log_a\!\left(\dfrac{x}{y}\right) = \log_a(x) - \log_a(y)$.
    \item \textbf{Logaritmo de una potencia:} $\log_a(x^n) = n\,\log_a(x)$.
    \item \textbf{Logaritmo de la base:} $\log_a(a) = 1$.
    \item \textbf{Logaritmo de 1:} $\log_a(1)=0$.
    \item \textbf{Cambio de base:} $\log_a(x) = \dfrac{\log_b(x)}{\log_b(a)}$, para cualquier $b>0$, $b\neq 1$.
    \item \textbf{Inversa:} $a^{\log_a(x)} = x$ y $\log_a(a^x)=x$.
    \item \textbf{Logaritmo natural y decimal:}
    \begin{itemize}
        \item $\ln(x) = \log_e(x)$, donde $e\approx 2.71828$.
        \item $\log(x)$ usualmente denota $\log_{10}(x)$.
    \end{itemize}
\end{enumerate}
\end{tcolorbox}

\begin{funcbox}{Ejemplo de aplicación}
\begin{align*}
    \log_2(8x) &= \log_2(8) + \log_2(x) = 3 + \log_2(x) \\
    \ln\!\left(\frac{e^3}{y}\right) &= \ln(e^3) - \ln(y) = 3 - \ln(y) \\
    \log_5(125) &= \log_5(5^3) = 3
\end{align*}
\end{funcbox}

%=====================================================================
% APARTADO 3: FUNCIÓN COMPUESTA
%=====================================================================
\section{Función Compuesta}

Dadas dos funciones $f$ y $g$, la función compuesta $(f\circ g)(x)=f(g(x))$ consiste en aplicar primero $g$ y luego $f$.

\subsection{Definición y dominio}
\begin{tcolorbox}[colback=orange!5, colframe=orange!70!black, title=Dominio de la composición]
Para que $(f\circ g)(x)$ esté definida, se requiere:
\begin{itemize}
    \item $x\in\Dominio(g)$.
    \item $g(x)\in\Dominio(f)$.
\end{itemize}
Es decir, $\Dominio(f\circ g) = \{x\in\Dominio(g) \mid g(x)\in\Dominio(f)\}$.
\end{tcolorbox}

\subsection{Ejemplos detallados}

\subsubsection{Ejemplo 1: $f(x)=\sqrt{x}$ y $g(x)=2x-3$}
\begin{align*}
    (f\circ g)(x) &= f(g(x)) = \sqrt{2x-3}.
\end{align*}
\begin{itemize}
    \item Dominio de $g$: $\mathbb{R}$.
    \item $g(x)\in\Dominio(f)=[0,\infty) \;\Rightarrow\; 2x-3 \geq 0 \;\Rightarrow\; x \geq \frac{3}{2}$.
    \item Por tanto, $\Dominio(f\circ g) = \left[\frac{3}{2},\infty\right)$.
    \item Recorrido: Como la raíz cuadrada devuelve valores $\geq 0$, el recorrido es $[0,\infty)$.
\end{itemize}

\subsubsection{Ejemplo 2: $f(x)=\ln(x)$ y $g(x)=x^2+1$}
\begin{align*}
    (f\circ g)(x) &= \ln(x^2+1).
\end{align*}
\begin{itemize}
    \item Dominio de $g$: $\mathbb{R}$.
    \item $g(x)=x^2+1 \geq 1 > 0$ para todo $x$, así que siempre pertenece al dominio de $\ln$, $(0,\infty)$.
    \item Luego $\Dominio(f\circ g)=\mathbb{R}$.
    \item Recorrido: como $x^2+1 \in [1,\infty)$, $\ln$ toma valores en $[\ln(1),\infty)=[0,\infty)$. Así $\Recorrido(f\circ g)=[0,\infty)$.
\end{itemize}

\subsubsection{Ejemplo 3: ¡Cuidado con el orden!}
Sean $f(x)=x^2$ y $g(x)=x+1$.
\begin{align*}
    (f\circ g)(x) &= (x+1)^2 = x^2+2x+1, \\
    (g\circ f)(x) &= x^2+1 \neq (f\circ g)(x).
\end{align*}
La composición no es conmutativa.

%=====================================================================
% APARTADO 4: FUNCIONES BIYECTIVAS Y RESTRICCIÓN
%=====================================================================
\section{Funciones Biyectivas y Restricción del Dominio}

\subsection{Conceptos previos}
Una función $f:A\to B$ es:
\begin{itemize}
    \item \textbf{Inyectiva} (uno a uno): si $f(x_1)=f(x_2) \Rightarrow x_1=x_2$. Gráficamente, cada recta horizontal corta la gráfica a lo sumo una vez.
    \item \textbf{Suprayectiva} (sobre): si $\Recorrido(f)=B$, es decir, todo elemento de $B$ es imagen de algún $x\in A$.
    \item \textbf{Biyectiva}: si es inyectiva y suprayectiva a la vez. Solo las funciones biyectivas admiten función inversa en todo su dominio.
\end{itemize}

\subsection{Restricción del dominio para lograr biyectividad}

Muchas funciones no son biyectivas de forma natural, pero al restringir su dominio y codominio adecuadamente, se convierten en biyectivas.

\subsubsection{Función cuadrática}
$f(x)=x^2$ con $f:\mathbb{R}\to\mathbb{R}$ no es inyectiva (porque $f(-2)=f(2)=4$) ni suprayectiva (su recorrido es $[0,\infty)\neq\mathbb{R}$).

\textbf{Restricción:} $f:[0,\infty)\to[0,\infty)$ con $f(x)=x^2$ sí es biyectiva. Es inyectiva (la rama derecha es creciente) y suprayectiva (cubre todos los no negativos).

\subsubsection{Función raíz cuadrada}
$g(x)=\sqrt{x}$ con $g:[0,\infty)\to[0,\infty)$ es naturalmente biyectiva (creciente y su recorrido llena $[0,\infty)$).

\subsubsection{Función racional homográfica}
$h(x)=\dfrac{2x+1}{x-3}$, con dominio natural $\mathbb{R}\setminus\{3\}$. Para que sea biyectiva, se toma codominio $\mathbb{R}\setminus\{2\}$ (excluyendo la asíntota horizontal). Así $h:\mathbb{R}\setminus\{3\}\to\mathbb{R}\setminus\{2\}$ es biyectiva (estrictamente monótona por ramas y su imagen es todo el codominio).

\subsubsection{Función valor absoluto}
$k(x)=|x|$ no es inyectiva ($|2|=|-2|=2$). Se restringe, por ejemplo, $k:[0,\infty)\to[0,\infty)$ y entonces $k(x)=x$ es biyectiva (identidad en no negativos). O bien $k:(-\infty,0]\to[0,\infty)$ con $k(x)=-x$.

%=====================================================================
% APARTADO 5: FUNCIONES INVERSAS
%=====================================================================
\section{Funciones Inversas}

Dada una función biyectiva $f$, su inversa $f^{-1}$ cumple:
\[
f^{-1}(f(x)) = x \quad \text{y} \quad f(f^{-1}(y)) = y.
\]
La gráfica de $f^{-1}$ es la reflexión de la gráfica de $f$ respecto a la recta $y=x$.

\subsection{Ejemplos de funciones inversas (basadas en el Apartado 1)}

\subsubsection{Lineal: $f(x)=2x-1$}
\begin{align*}
    y &= 2x-1 \;\Rightarrow\; y+1 = 2x \;\Rightarrow\; x = \frac{y+1}{2}.
\end{align*}
Por tanto, $f^{-1}(x)=\dfrac{x+1}{2}$. Dominio y recorrido: ambos $\mathbb{R}$.

\begin{center}
\begin{tikzpicture}
    \begin{axis}[
        axis lines=middle,
        xlabel=$x$, ylabel=$y$,
        xmin=-3, xmax=3,
        ymin=-3, ymax=3,
        grid=both,
        minor tick num=1,
        width=7cm, height=7cm,
        title={$f$ y $f^{-1}$ lineales}
    ]
    \addplot[domain=-2.5:2.5, blue, thick] {2*x-1};
    \addplot[domain=-2.5:2.5, red, thick] {(x+1)/2};
    \addplot[domain=-3:3, dashed, gray] {x};
    \legend{$f(x)=2x-1$, $f^{-1}(x)=\frac{x+1}{2}$, $y=x$}
    \end{axis}
\end{tikzpicture}
\end{center}

\subsubsection{Cuadrática restringida: $f(x)=x^2$ en $[0,\infty)$}
\begin{align*}
    y &= x^2,\; x\geq 0 \;\Rightarrow\; x = \sqrt{y}.
\end{align*}
Inversa: $f^{-1}(x)=\sqrt{x}$ con dominio $[0,\infty)$.

\subsubsection{Exponencial y logarítmica}
$f(x)=e^x$ (dominio $\mathbb{R}$, recorrido $(0,\infty)$) es biyectiva. Su inversa es $f^{-1}(x)=\ln(x)$, dominio $(0,\infty)$, recorrido $\mathbb{R}$.

\begin{center}
\begin{tikzpicture}
    \begin{axis}[
        axis lines=middle,
        xlabel=$x$, ylabel=$y$,
        xmin=-3, xmax=5,
        ymin=-3, ymax=5,
        grid=both,
        minor tick num=1,
        width=7cm, height=7cm,
        title={Exponencial y logaritmo}
    ]
    \addplot[domain=-3:1.6, blue, thick] {exp(x)};
    \addplot[domain=0.05:5, red, thick] {ln(x)};
    \addplot[domain=-3:5, dashed, gray] {x};
    \legend{$e^x$, $\ln(x)$, $y=x$}
    \end{axis}
\end{tikzpicture}
\end{center}

\subsubsection{Racional: $f(x)=\dfrac{2x+1}{x-3}$}
Dominio $\mathbb{R}\setminus\{3\}$, recorrido $\mathbb{R}\setminus\{2\}$. Despejamos $x$:
\begin{align*}
    y(x-3) &= 2x+1 \;\Rightarrow\; yx - 3y = 2x+1 \\
    yx - 2x &= 3y+1 \;\Rightarrow\; x(y-2) = 3y+1 \\
    x &= \frac{3y+1}{y-2}.
\end{align*}
Inversa: $f^{-1}(x)=\dfrac{3x+1}{x-2}$ con dominio $\mathbb{R}\setminus\{2\}$.

\subsubsection{Raíz cúbica y cúbica}
$f(x)=x^3$ (biyectiva en $\mathbb{R}$) tiene inversa $f^{-1}(x)=\sqrt[3]{x}$, ambas con dominio y recorrido $\mathbb{R}$.

\subsubsection{Cúbica desplazada: $f(x)=(x-1)^3+2$}
\begin{align*}
    y &= (x-1)^3+2 \;\Rightarrow\; y-2 = (x-1)^3 \\
    \sqrt[3]{y-2} &= x-1 \;\Rightarrow\; x = \sqrt[3]{y-2}+1.
\end{align*}
Inversa: $f^{-1}(x)=\sqrt[3]{x-2}+1$.

\subsubsection{Valor absoluto restringido}
Si restringimos $f(x)=|x|$ a $x\geq 0$, la función es $f(x)=x$, cuya inversa es $f^{-1}(x)=x$ (es su propia inversa en ese dominio).

%=====================================================================
% NUEVA SECCIÓN: SECCIONES CÓNICAS
%=====================================================================
\newpage
\section{Secciones Cónicas}

Las cónicas son curvas obtenidas al cortar un cono con un plano. En su forma canónica con centro (o vértice) en $(h,k)$, sus ecuaciones son:

\subsection{Circunferencia}
\begin{funcbox}{Circunferencia $(x-h)^2+(y-k)^2=r^2$}
\begin{minipage}[t]{0.55\textwidth}
    \textbf{Ecuación:} $(x-h)^2+(y-k)^2=r^2$\\
    \textbf{Centro:} $(h,k)$\\
    \textbf{Radio:} $r>0$\\
    \textbf{Características:}
    \begin{itemize}
        \item Lugar geométrico de puntos a distancia $r$ del centro.
        \item Gráfica: curva cerrada.
    \end{itemize}
    \textbf{Ejemplo:} $(x-1)^2+(y+2)^2=9$ \\ Centro $(1,-2)$, radio $3$.
\end{minipage}\hfill
\begin{minipage}[t]{0.4\textwidth}
    \centering
    \begin{tikzpicture}
        \begin{axis}[
            axis lines=middle,
            xlabel=$x$, ylabel=$y$,
            xmin=-3, xmax=5,
            ymin=-6, ymax=2,
            grid=both,
            minor tick num=1,
            width=6cm, height=5cm,
            title={$(x-1)^2+(y+2)^2=9$}
        ]
        \addplot[domain=0:360, samples=200, blue, thick] 
            ({1+3*cos(x)},{-2+3*sin(x)});
        \addplot[mark=*, only marks, red] coordinates {(1,-2)};
        \end{axis}
    \end{tikzpicture}
\end{minipage}
\end{funcbox}

\subsection{Parábola}
\subsubsection{Parábola vertical}
\begin{funcbox}{Parábola vertical $(x-h)^2=4p(y-k)$}
\begin{minipage}[t]{0.55\textwidth}
    \textbf{Ecuación:} $(x-h)^2=4p(y-k)$ o $y=a(x-h)^2+k$ con $a=\frac{1}{4p}$.\\
    \textbf{Vértice:} $(h,k)$\\
    \textbf{Foco:} $(h, k+p)$ (si abre hacia arriba)\\
    \textbf{Directriz:} $y=k-p$\\
    \textbf{Orientación:} $p>0$ abre hacia arriba; $p<0$ hacia abajo.\\
    \textbf{Ejemplo:} $(x-2)^2=8(y-1)$; $p=2$, vértice $(2,1)$, foco $(2,3)$.
\end{minipage}\hfill
\begin{minipage}[t]{0.4\textwidth}
    \centering
    \begin{tikzpicture}
        \begin{axis}[
            axis lines=middle,
            xlabel=$x$, ylabel=$y$,
            xmin=-2, xmax=6,
            ymin=-2, ymax=6,
            grid=both,
            minor tick num=1,
            width=6cm, height=5cm,
            title={$(x-2)^2=8(y-1)$}
        ]
        \addplot[domain=-0.8:4.8, samples=200, blue, thick] 
            {1 + (x-2)^2/8};
        \addplot[mark=*, only marks, red] coordinates {(2,1)};
        \end{axis}
    \end{tikzpicture}
\end{minipage}
\end{funcbox}

\subsubsection{Parábola horizontal}
\begin{funcbox}{Parábola horizontal $(y-k)^2=4p(x-h)$}
\begin{minipage}[t]{0.55\textwidth}
    \textbf{Ecuación:} $(y-k)^2=4p(x-h)$\\
    \textbf{Vértice:} $(h,k)$\\
    \textbf{Foco:} $(h+p,k)$ (si abre a la derecha)\\
    \textbf{Directriz:} $x=h-p$\\
    \textbf{Ejemplo:} $(y+1)^2=12(x-3)$; $p=3$, vértice $(3,-1)$, abre a la derecha.
\end{minipage}\hfill
\begin{minipage}[t]{0.4\textwidth}
    \centering
    \begin{tikzpicture}
        \begin{axis}[
            axis lines=middle,
            xlabel=$x$, ylabel=$y$,
            xmin=-1, xmax=9,
            ymin=-8, ymax=6,
            grid=both,
            minor tick num=1,
            width=6cm, height=5cm,
            title={$(y+1)^2=12(x-3)$}
        ]
        \addplot[
            domain=-7:5, 
            samples=200, 
            smooth, 
            blue, 
            thick, 
            variable=\t
        ]
        ({3+(\t+1)^2/12},{\t}); % parametrización correcta
        \end{axis}
    \end{tikzpicture}
\end{minipage}
\end{funcbox}

\subsection{Elipse}
\begin{funcbox}{Elipse $\dfrac{(x-h)^2}{a^2}+\dfrac{(y-k)^2}{b^2}=1$}
\begin{minipage}[t]{0.55\textwidth}
    \textbf{Ecuación:} $\dfrac{(x-h)^2}{a^2}+\dfrac{(y-k)^2}{b^2}=1$\\
    \textbf{Centro:} $(h,k)$\\
    \textbf{Semieje mayor:} $a$, semieje menor: $b$ (si $a>b$ es horizontal; si $b>a$ vertical).\\
    \textbf{Ejemplo:} $\dfrac{(x-2)^2}{16}+\dfrac{(y+1)^2}{9}=1$\\
    Centro $(2,-1)$, semieje mayor horizontal $a=4$, menor $b=3$.
\end{minipage}\hfill
\begin{minipage}[t]{0.4\textwidth}
    \centering
    \begin{tikzpicture}
        \begin{axis}[
            axis lines=middle,
            xlabel=$x$, ylabel=$y$,
            xmin=-3, xmax=7,
            ymin=-5, ymax=3,
            grid=both,
            minor tick num=1,
            width=6cm, height=5cm,
            title={Elipse $(x-2)^2/16+(y+1)^2/9=1$}
        ]
        \addplot[domain=0:360, samples=200, blue, thick] 
            ({2+4*cos(x)},{-1+3*sin(x)});
        \addplot[mark=*, only marks, red] coordinates {(2,-1)};
        \end{axis}
    \end{tikzpicture}
\end{minipage}
\end{funcbox}

\subsection{Hipérbola}
\subsubsection{Hipérbola horizontal}
\begin{funcbox}{Hipérbola $\dfrac{(x-h)^2}{a^2}-\dfrac{(y-k)^2}{b^2}=1$}
\begin{minipage}[t]{0.55\textwidth}
    \textbf{Ecuación:} $\dfrac{(x-h)^2}{a^2}-\dfrac{(y-k)^2}{b^2}=1$\\
    \textbf{Centro:} $(h,k)$\\
    \textbf{Vértices:} $(h\pm a, k)$\\
    \textbf{Asíntotas:} rectas que pasan por el centro con pendiente $\pm b/a$.\\
    \textbf{Ejemplo:} $\dfrac{(x-1)^2}{9}-\dfrac{(y-2)^2}{4}=1$\\
    Centro $(1,2)$, $a=3$, $b=2$.
\end{minipage}\hfill
\begin{minipage}[t]{0.4\textwidth}
    \centering
    \begin{tikzpicture}
        \begin{axis}[
            axis lines=middle,
            xlabel=$x$, ylabel=$y$,
            xmin=-5, xmax=7,
            ymin=-4, ymax=8,
            grid=both,
            minor tick num=1,
            width=6cm, height=5cm,
            title={$\frac{(x-1)^2}{9}-\frac{(y-2)^2}{4}=1$}
        ]
        % Rama derecha
        \addplot[domain=-1:1, samples=200, blue, thick] 
            ({1+3*cosh(x)},{2+2*sinh(x)});
        % Rama izquierda
        \addplot[domain=-1:1, samples=200, blue, thick] 
            ({1-3*cosh(x)},{2+2*sinh(x)});
        % Asíntotas
        \addplot[dashed, domain=-4:6, samples=2] {2 + 2/3*(x-1)};
        \addplot[dashed, domain=-4:6, samples=2] {2 - 2/3*(x-1)};
        \end{axis}
    \end{tikzpicture}
\end{minipage}
\end{funcbox}

\subsubsection{Hipérbola vertical}
\begin{funcbox}{Hipérbola $\dfrac{(y-k)^2}{a^2}-\dfrac{(x-h)^2}{b^2}=1$}
\begin{minipage}[t]{0.55\textwidth}
    \textbf{Ecuación:} $\dfrac{(y-k)^2}{a^2}-\dfrac{(x-h)^2}{b^2}=1$\\
    \textbf{Centro:} $(h,k)$\\
    \textbf{Vértices:} $(h, k\pm a)$\\
    \textbf{Asíntotas:} $y-k = \pm \frac{a}{b}(x-h)$.\\
    \textbf{Ejemplo:} $\dfrac{(y+1)^2}{4}-\dfrac{(x-3)^2}{1}=1$\\
    Centro $(3,-1)$, $a=2$, $b=1$, abre hacia arriba y abajo.
\end{minipage}\hfill
\begin{minipage}[t]{0.4\textwidth}
    \centering
    \begin{tikzpicture}
        \begin{axis}[
            axis lines=middle,
            xlabel=$x$, ylabel=$y$,
            xmin=-1, xmax=7,
            ymin=-6, ymax=4,
            grid=both,
            minor tick num=1,
            width=6cm, height=5cm,
            title={$\frac{(y+1)^2}{4}-\frac{(x-3)^2}{1}=1$}
        ]
        % Rama superior
        \addplot[domain=-1:1, samples=200, blue, thick] 
            ({3+sinh(x)},{-1+2*cosh(x)});
        % Rama inferior
        \addplot[domain=-1:1, samples=200, blue, thick] 
            ({3+sinh(x)},{-1-2*cosh(x)});
        % Asíntotas
        \addplot[dashed, domain=-0.5:6.5, samples=2] {-1 + 2*(x-3)};
        \addplot[dashed, domain=-0.5:6.5, samples=2] {-1 - 2*(x-3)};
        \end{axis}
    \end{tikzpicture}
\end{minipage}
\end{funcbox}

\end{document}